\section{Reproducibility, proof status, and comparison}\label{sec:provenance}

\subsection{Independent computations}
The archive contains complete \LaTeX\ sources and compiled PDFs, an exact rational
law calculator, independent finite-field enumeration, recorded results, and a
formalization plan. The calculator groups labels by sufficient statistics.
The independent checker instead constructs actual vector assignments and counts
their edge observations; it does not use a proposed formula to generate those
counts. It includes the extension field
$\F_4=\F_2[z]/(z^2+z+1)$.

The direct enumeration verified:
\begin{itemize}
\item $1{,}745{,}891$ nonzero-vector assignments across $16$ cycle cases, four
cactus cases, and five tree cases;
\item all $1{,}976$ resulting cycle label probabilities, their exact total
variation and chi-square divergences;
\item $608$ proper cycle-edge subsets, giving $11{,}614$ marginal equalities;
\item six additive-character inversions for even-cycle resonance.
\end{itemize}
The fields were $\F_2,\F_3,\F_4,\F_5$, and ambient dimensions were two or four.
All comparisons passed. A separate $432$-case arithmetic check compared the
grouped likelihoods with the closed formulas for
$q\in\{2,3,4,5,7,8,9,16,25\}$ and several dimensions and cycle lengths.

A further twenty complete-bipartite cases enumerate $4{,}094{,}368$ assignments and check $7{,}382$ full label probabilities. Thus the two graph verifiers together cover $5{,}840{,}259$ actual assignments. They also check sixteen limiting-score variances by exact summation over $5{,}680$ label vectors.

These finite checks are regression evidence, not the proof for arbitrary
$q,n$. The proof is the character calculation, exact finite summation,
block factorization, and probability-limit argument. No Monte Carlo
approximation is used to establish an exact identity or the normal limit.
Further companion checks, where included, are reported separately in their
own files so their scope is not conflated with these totals.

\subsection{Source audit and novelty scope}
The pinned repository distinguishes unconditional results, uses of Dirichlet
GRH, and reductions whose target estimates remain open. Its README describes
the large article as a merge of reports, not an independent verification of
every claim \cite{ProveItReadme}. The vendored upstream is pinned at
\texttt{72fead102f663d5c7a407964a13dc0cf9d977207}; this work did not rerun that
development's kernel checks. Our new theorems do not require them.

The principal mathematical input from the repository is a method: evaluate
finite-field Fourier transforms using bilinear orthogonality, then retain
the structured information needed to compute the full law. The identity
$F_sF_t=NP_{s/t}-J$ makes this transfer exact on graph cycles. We prove it
directly. Consequently a revision elsewhere in the analytic corpus would not
affect these results.

Related primary literature includes cycle counts for distance and dot-product
graphs \cite{IJM}, extensions to general nondegenerate bilinear forms
\cite{Bright}, classical finite symplectic graph theory \cite{TangWan}, and
topology-dependent distinguishability for real masked Wishart matrices
\cite{BBH}. Those works establish relevant background. Special fixed-label cycle
counts can also follow from known graph spectra. The contribution here is the
complete labelled finite-field law relative to the product of its true
marginals, followed by the exact critical profile and explicit overlapping-cycle
contrast. We did not find this combined statement in the inspected repository
or the specifically reviewed references. This is a scoped novelty assessment,
not a proof of worldwide bibliographic priority.

The results have not been encoded and checked in Lean. The formalization plan
separates finite algebra from probability limits and states what needs formal
proof. It contains no placeholder presented as a verified theorem.

\subsection{Scope and remaining questions}
The cactus hypothesis is substantive. Our complete-bipartite calculation handles
one family of overlapping cycles, but does not give a general answer for every
graph. The bounded cycle-length hypothesis controls all likelihood increments
uniformly. If the lengths grow with $n$, further estimates are needed.
Nonuniform vector laws, zero vectors, projective sampling, and distinctness
constraints on graph-vertex vectors are different models.

A natural next problem is to determine the information contributed by overlapping
cycles in more general graphs. The complete-bipartite example shows that adding
the individual cycle chi-square values can fail drastically. The unconditional
graph likelihood can have extremely rare, very large values even when its total
variation distance tends to zero. Any wider theorem must account for this effect.

The archive also includes a separate companion theorem about the full
cross-pairing matrix of arbitrarily many prescribed small frames. Its assumptions,
proof, comparison with earlier rank-universality literature, and numerical
checks are self-contained in that paper.

